\documentclass[12pt]{article}
\usepackage[utf8]{inputenc}
\usepackage[T1]{fontenc}
\usepackage[margin=1in]{geometry}
\usepackage{setspace}
\usepackage{xcolor}
\usepackage[colorlinks=true, urlcolor=blue!60!black, linkcolor=blue!60!black]{hyperref}
\setlength{\parindent}{2em}
\setlength{\parskip}{1.25em}
\renewcommand{\baselinestretch}{1.0}
\begin{document}
\begin{center}
{\LARGE \textit{Hans v. Louisiana}, 134 U.S. 1 (1890)\par}
\vspace{0.8em}
{\large Case Brief --- Four-Page Edition\par}
\end{center}
\vspace{1em}

\section{Memory Jogger}

Despite the Eleventh Amendment's text referencing only suits by out-of-state or foreign citizens, the background constitutional principle of state sovereign immunity bars a citizen from suing his own state in federal court without its consent, even on a federal-question claim (\href{https://supreme.justia.com/cases/federal/us/134/1/}{134 U.S. 1}).

\section{Detailed Case Facts}

Under an 1874 act, Louisiana issued ``consolidated bonds'' payable in forty years with semiannual interest coupons, and a companion constitutional amendment declared the bonds a valid, unimpeachable contract and dedicated an annual tax to pay them. After Reconstruction, Louisiana's 1879 constitution remitted the coupons due January 1, 1880, and diverted the taxes already collected for them to general state expenses.

Bernard Hans, a Louisiana citizen holding those coupons, sued the State in December 1884 for \$87,500 plus interest. He alleged that the 1879 repudiation impaired the obligation of his contract in violation of the Contract Clause and invoked federal-question jurisdiction under Article III and the Judiciary Act of 1875.

\section{Procedural History}

Hans sued in the federal circuit court for the Eastern District of Louisiana. The state Attorney General excepted that the State could not be sued without its consent, and in 1885 Judge Billings sustained the exception and dismissed the suit. On writ of error, the Supreme Court heard argument on January 22, 1890, and unanimously affirmed on March 3, 1890.

\section{Judicial Votes}

Affirmed, \textbf{9--0} (8--1 on reasoning).
\begin{itemize}
\item \textbf{Majority (8):} Justice Bradley, joined by Chief Justice Fuller and Justices Miller, Field, Gray, Blatchford, Lamar, and Brewer.
\item \textbf{Concurrence in the judgment:} Justice Harlan.
\item \textbf{Dissents:} None.
\end{itemize}

\section{Holding}

An unconsenting State cannot be sued in federal court by one of its own citizens, even when the suit arises under the Constitution or laws of the United States. Neither Article III's federal-question grant nor the Judiciary Act of 1875 authorizes such suits, because the Constitution presupposes the States' preexisting sovereign immunity, an attribute not surrendered in the plan of the Convention. Bradley's disapproval of \textit{Chisholm v. Georgia} and his remark that violations of contract rights ``may be judicially resisted'' are dicta.

\section{Analysis of Opinions}

\textbf{Majority (Bradley).} Bradley acknowledged that the Eleventh Amendment literally bars only suits by ``Citizens of another State, or by Citizens or Subjects of any Foreign State,'' but rejected a literal reading as producing an ``anomalous result'': a State could be sued by its own citizens on claims it could not be sued on by outsiders, and in courts where it did not permit itself to be sued at home. \textit{Chisholm} had caused ``such a shock of surprise throughout the country'' that the Amendment was swiftly adopted, declaring that the Constitution ``shall not be construed'' to permit private suits against States and vindicating Iredell's dissent. Bradley drew on Hamilton's \textit{Federalist No.\ 81} and on assurances by Madison and John Marshall at the Virginia convention, and called a proviso allowing suits by a State's own citizens ``almost an absurdity on its face.''

He added three points. Precedent confirmed that unconsented suits against a State were unknown. The 1875 Act gave jurisdiction ``concurrent with the courts of the several States,'' which could not hear such suits. And Marshall's contrary suggestion in \textit{Cohens v. Virginia} was ``unnecessary to the decision, and in that sense extrajudicial.'' A State's contracts rest on its political honor; courts may not compel it to pay its debts.

\textbf{Harlan (concurring in the judgment).} Harlan agreed only that federal judicial power does not reach a citizen's direct suit against his unconsenting State. He refused to join ``many things said in the opinion,'' because the attack on \textit{Chisholm} was unnecessary and \textit{Chisholm} ``was based upon a sound interpretation of the Constitution as that instrument then was.''

\section{Examples: Future Applications}

\begin{enumerate}
\item \textbf{Same-Side (Barred): Overtime Damages.} A Texas resident sues Texas in federal court for retrospective damages under a federal overtime statute enacted under the Commerce Clause. Under \textit{Hans}, reinforced by \textit{Seminole Tribe} (1996), immunity bars the suit and Congress cannot abrogate it under Article I.
\item \textbf{Same-Side (Barred): Pension Back Pay.} A state employee sues the Louisiana Department of Revenue, an arm of the State, for back pay on a Contract Clause theory, seeking money from the treasury.
\item \textbf{Opposite-Side (Permitted): Officer Injunction.} A citizen seeks only a prospective injunction against the Attorney General's enforcement of an unconstitutional law. \textit{Hans} distinguished suits against officers, and \textit{Ex parte Young} (1908) permits them.
\item \textbf{Opposite-Side (Permitted): Counties and the United States.} A creditor sues a county over defaulted bonds, or the United States sues a State. Immunity does not protect political subdivisions (\textit{Lincoln County v. Luning}, decided the same day), and States surrendered immunity to federal suits.
\item \textbf{Fence-Sitter: Plan-of-the-Convention Statutes.} Congress authorizes private suits against States under war or eminent-domain powers. \textit{Hans} ordinarily forbids this, but immunity surrendered ``in the plan of the convention'' has supported divided rulings in \textit{Katz}, \textit{PennEast}, and \textit{Torres}.
\end{enumerate}

\section{Critique}

\textbf{Originalist and textualist views.} Diversity theorists such as William Fletcher, John Gibbons, and Akhil Amar argue that \textit{Hans} contradicts the text: the Amendment repealed only \textit{Chisholm}'s diversity jurisdiction, leaving federal-question suits untouched. Original-law scholars William Baude, Stephen Sachs, and Caleb Nelson defend the result as resting on a background common-law immunity, though that view may give Congress more power to abrogate than \textit{Seminole Tribe} and \textit{Alden} allowed.

\textbf{Progressive and doctrinal views.} Vicki Jackson, Edward Purcell, and John Orth see \textit{Hans} as a post-Reconstruction retreat, after Redeemer governments repudiated bonds the Court could not force them to pay. Elevating it into a constitutional shield, critics argue, undermines the Supremacy Clause and \textit{Marbury}'s remedial maxim.

\section{Key Quotations}

\begin{quote}
``The suability of a State without its consent was a thing unknown to the law.''
\end{quote}
\noindent --- Justice Bradley, Opinion of the Court, 134 U.S. at 16.

\begin{quote}
``Can we suppose that, when the Eleventh Amendment was adopted, it was understood to be left open for citizens of a State to sue their own state in the federal courts, whilst the idea of suits by citizens of other states, or of foreign states, was indignantly repelled? Suppose that Congress, when proposing the Eleventh Amendment, had appended to it a proviso that nothing therein contained should prevent a State from being sued by its own citizens in cases arising under the Constitution or laws of the United States, can we imagine that it would have been adopted by the States? The supposition that it would is almost an absurdity on its face.''
\end{quote}
\noindent --- Justice Bradley, Opinion of the Court, 134 U.S. at 15.

\begin{quote}
``While the State cannot be compelled by suit to perform its contracts, any attempt on its part to violate property or rights acquired under its contracts may be judicially resisted, and any law impairing the obligation of contracts under which such property or rights are held is void and powerless to affect their enjoyment.''
\end{quote}
\noindent --- Justice Bradley, Opinion of the Court, 134 U.S. at 20--21.

\end{document}
